\section{The finite-real question}

A periodic-orbit product can change when one orbit is removed, even if the
underlying measure space is unchanged modulo null sets. Generalized Boole maps
provide an explicit setting in which this distinction persists through a
parabolic transition. For $a,b>0$ set
\begin{equation}\label{eq:domain}
 T_{a,b}(x)=ax-\frac b x,
 \qquad X_{a,b}=\R\setminus\bigcup_{j\ge0}T_{a,b}^{-j}\{0\}.
\end{equation}
Inverse images in this definition are computed for the rational map on the
sphere, and only finite real points are deleted. We study its restriction to
$X_{a,b}$, with one unit of time equal to one original iteration. Infinity is
not a physical state. Write $q=2a-1$ and define
\begin{equation}\label{eq:definition}
 \tau_n(a)=\sum_{x\in\Fix_{X_{a,b}} T_{a,b}^n}
 \frac{1}{(T_{a,b}^n)'(x)-1},\qquad
 D_a(u)=\exp\left(-\sum_{n\ge1}\frac{\tau_n(a)}n u^n\right).
\end{equation}
We prove below that every denominator is positive and that the series converges
on $|u|<1$. This notation defines a stability germ, not an operator determinant
by assumption. The scale $b$ disappears from the formulas because
$x=\sqrt b\,y$ conjugates $T_{a,b}$ to $T_{a,1}$ without changing time.

\begin{theorem}[Stability weights and continuation]\label{thm:main}
For all $a,b>0$ and every $n\ge1$, every finite real fixed point of $T_{a,b}^n$
is simple and has multiplier greater than one. Its weighted sum is
\begin{equation}\label{eq:tau-main}
 \tau_n(a)=
 \begin{cases}
 \displaystyle\frac2{1-q^n}-\frac1{1-a^n},&0<a<1,\\[4pt]
 \displaystyle\frac{n-1}{2n},&a=1,\\[4pt]
 \displaystyle\frac1{a^n-1},&a>1.
 \end{cases}
\end{equation}
The normalized germ has the following continuations:
\begin{equation}\label{eq:products-main}
 D_a(u)=
 \begin{cases}
 \displaystyle(1-u)\frac{\prod_{j\ge1}(1-q^ju)^2}
 {\prod_{k\ge1}(1-a^ku)},&0<a<1,\\[7pt]
 \displaystyle(1-u)^{1/2}\exp(\Li(u)/2),&a=1,\\[4pt]
 \displaystyle\prod_{k\ge1}(1-a^{-k}u),&a>1.
 \end{cases}
\end{equation}
For $0<a<1$, the product is entire if and only if
\begin{equation}\label{eq:entire-main}
 0<a<\frac12,\qquad a=(1-2a)^{2m}\quad\hbox{for some }m\in\N.
\end{equation}
For each $m\ge1$ there is exactly one such parameter; the first is $a=1/4$.
All remaining zero and pole multiplicities are classified in
Proposition~\ref{prop:divisor}.
\end{theorem}

\begin{theorem}[Compensated parabolic limit]\label{thm:critical}
The critical germ $D_1$ is not meromorphic through $u=1$. For $a>1$ define
\begin{equation}\label{eq:reduced-definition}
 F_a(u)=\prod_{j\ge1}(1-q^{-j}u)^2,\qquad
 D_a^{\red}(u)=\frac{D_a(u)}{F_a(u)}.
\end{equation}
The factor $F_a$ is exactly the complete primitive stability factor of the two
finite real fixed orbits, including all repetitions. On $|u|<1$,
\begin{equation}\label{eq:limits-main}
 D_a\longrightarrow D_1\quad(a\uparrow1),\qquad
 D_a^{\red}\longrightarrow D_1\quad(a\downarrow1)
\end{equation}
locally uniformly. The unreduced supercritical germs do not have this limit.
No statement of uniform convergence includes the boundary point $u=1$.
\end{theorem}

\paragraph{Classical inputs and the scope of the increment.}
The statistical phase picture is established prior work: Umeno and Okubo
describe the Cauchy, infinite-measure and dissipative regimes and compute
Lyapunov exponents~\citep{umeno2016boole}. Mendoza Mendoza and Ruiz Leal study
the more general map $ax-b/x+c$~\citep{mendoza2022boole}. Their deleted-prepole
domain (Section~3) and symbolic descriptions are directly relevant here. In
particular, their Theorems~4.2 and~4.4 for the critical and $1/2<a<1$ regimes
already yield the corresponding unweighted periodic counts; their Theorem~4.5
gives the supercritical binary Cantor dynamics. We do not claim the domain, phase structure,
symbolic coding or those counting consequences as new results.

The closest operator input is the finite-Blaschke spectrum of Bandtlow, Just
and Slipantschuk~\citep{bandtlow2017spectral}. It concerns natural holomorphic
function spaces on the full circle. In the subcritical regime the Cayley map
conjugates $T_{a,b}$ to a Blaschke product, but the physical periodic set omits
the circle point corresponding to infinity. Section~\ref{sec:scope} gives its
exact missing factor. This deletion must not be inferred from an $L^p$ quotient:
removing a null orbit alone does not remove a trace contribution.

The single question addressed here is the complete \emph{weighted} finite-real
product across the transition: the exact divisor, its exceptional entire
parameters, and the two-sided limit with the dynamically specified
supercritical compensation. We derive these statements without invoking a
natural transfer-operator realization on $X_{a,b}$. The literature comparison
is bounded and does not establish publication priority. Sections~2--3 identify
the physical cycles and compute their indices; Sections~4--5 prove the
continuation and limiting statements.
